\documentclass[11pt,a4paper]{article}
\usepackage[utf8]{inputenc}
\usepackage[T1]{fontenc}
\usepackage[spanish]{babel}
\usepackage{hyperref}
\usepackage{geometry}
\geometry{margin=2.5cm}
\usepackage{parskip}

\title{084 FARO- Fuente, Acto, Rastro, Orden}
\author{Alejandro Valenzuela Rivera\\\small Universidad Autónoma de Nuevo León}
\date{\today}
\begin{document}
\maketitle

\begin{quote}
\textbf{Sesgo declarado.} Este trabajo se sostiene desde el sesgo declarado de su autor, y no finge neutralidad. Lo que afirma se puede tumbar, y ese es el punto: el falsificador está escrito, no escondido.
\end{quote}

\section*{Resumen}
El nombre del lenguaje de procesos. No dice lo que todavia no esta: dice como se llega. Fuente es de donde viene cada pieza; Acto es lo que se ejecuta; Rastro es donde quedo y quien lo lee; Orden es que la consume y en que borde.

\section{Por qué se sostiene}
Un nombre en ingles no es de aqui, y el acronimo tiene que decir el acto en el idioma en que se trabaja. FARO es corto, se pronuncia de un tiron, y dice a donde mirar: la fuente de la que viene, el acto que se hace, el rastro que deja y el orden que sigue.

\section{De dónde viene}
La autoridad, y el repo.

\section{El caso}
Ninguna: es el nombre de la carpeta.

\section{Método}
Una biblioteca que dice de donde viene cada pieza, y que se entrega.

\section{Límites, y qué lo tumba}
Una biblioteca sin version, que no se puede citar sin ambiguedad.

\section{Lo que queda abierto}
\textquestiondown{}A donde mira este lenguaje: que fuente, que acto, que rastro y que orden?

\section*{Cómo se leyó esta tesis}
Esto se escribió en el repositorio man-made-empathy, del que este candado es una pieza: candados/06 del caso/084 FARO- Fuente, Acto, Rastro, Orden.md. La máquina que lo compuso lo hizo desde este mismo archivo, y no agregó nada que no estuviera escrito. Si algo de esto es falso, se dice en el mismo sitio, y se corrige con fecha.

\end{document}
